\section{Lecture 11}

Performance of SGD at small \(\gamma\)
\[\theta_i = \theta_{i-1} - \gamma g(\theta_{i-1}, x_i)\]
Asusme that \(m(\theta) = \E[g(\theta, X_i)]\) is \(\mu\)-strongly convex and that \(m(\theta^*) = 0\) and \(\theta^*= 0\). Furthermore, assume
\[\|g(\theta, X_i) - g(\theta', X_i)\|_p \leq L_p |\theta - \theta'|\]
Then \(\theta_i(\gamma)\) converges to a stationary distribution \(\pi_\gamma\). We want to ascertain whether
\[\frac{1}{\sqrt \gamma}\theta_\infty(\gamma) \Rightarrow \N(0, \Sigma)\]
as \(\gamma \rightarrow 0\). A special case occurs when
\[\begin{gathered}\theta_i^L = \theta_{i-1}^L - \gamma (\theta_{i-1}^L - X_i) \\
= \gamma X_i + (1 - \gamma)\gamma X_{i-1} + (1 - \gamma)^2 \gamma X_{i-2} + \cdots\end{gathered}\]
It\textquotesingle s easy to prove a CLT in this setting. We have
\[\begin{gathered}\Var[\theta_i^L] = \gamma^2 \sigma^2 + (1 - \gamma)^2 \gamma^2 \sigma^2 + (1 - \gamma)^4 \gamma^2 \sigma^2 + \cdots \\
= \frac{\gamma^2 \sigma^2}{1 - (1 - \gamma)^2} \\
\approx \frac{\gamma \sigma^2}{2 - \gamma} \\
\approx \frac{\gamma \sigma^2}{2}.\end{gathered}\]
Thus,
\[\frac{\theta_i^L}{\sqrt \gamma} \Rightarrow \N(0, \sigma^2/2)\]
This follows immediately from the Lindeberg-Feller CLT. Next, we need to linearize to tackle the general case. We can rewrite as
\[\theta_i = \theta_{i-1} - \gamma m(\theta_{i-1}) + \gamma [m(\theta_{i-1}) - g(\theta_{i-1}, X_i)]\]
Next, recall Rio\textquotesingle s inequality: \(\|\theta_i\|_p^2 = O(\gamma)\). Then
\[= \theta_{i-1} - \gamma m(\theta_{i-1}) + \gamma[m(0) - g(0, X_i)] + \gamma D(\theta_{i-1}, X_i)\]
Now,
\[\|D(\theta_{i-1}, X_i)\|_p \leq 2\| L_p |\theta_{i-1} - 0|\|_p = 2L_p\|\theta_{i-1}\|_p\]
Call \(\eta_i := m(0) - g(0, X_i)\). Then
\[\theta_i^\diamond = \theta_{i-1}^\diamond - \gamma m(\theta_{i-1}^\diamond) + \gamma \eta_i\]
Note that \(\eta_i\) is an independent sequence. We call \(\theta_i^\diamond\) a nonlinear AR(1) process. Next,
\[\|\theta_i - \theta_i^\diamond\|_p^2 = \|(\theta_{i-1} - \gamma m(\theta_{i-1})) + \gamma D(\theta_{i-1}, X_i) - (\theta_{i-1}^\diamond - \gamma m(\theta_{i-1}^\diamond))\|_p^2\]
By Rio\textquotesingle s inequality,
\[\leq \| (\theta_{i-1} - \gamma m(\theta_{i-1})) - (\theta_{i-1}^\diamond - \gamma m(\theta_{i-1}^\diamond))\|_p^2 + (p - 1)\|\gamma D(\theta_{i-1}, X_i)\|_p^2\]
Using convexity,
\[\| (\theta_{i-1} - \gamma m(\theta_{i-1})) - (\theta_{i-1}^\diamond - \gamma m(\theta_{i-1}^\diamond))\|_p^2 \leq \|(1 - \gamma \kappa)(\theta_{i-1} - \theta_{i-1}^\diamond)\|_p^2\]
Since \(\|D(\theta_{i-1}, X_i)\|_p \leq \gamma\), we have
\[(p - 1)\|\gamma D(\theta_{i-1}, X_i)\|_p^2 \leq (p - 1)\gamma^2 \gamma\]
Let \(\delta_i = \|\theta_i - \theta_i^\diamond|_p\). Then we have derived
\[\delta_i \leq (1 - \gamma \kappa)^2 \delta_{i-1} + (p - 1)O(\gamma^3)\]
This shows that \(\delta_i = O(\gamma)\) and thus justifies replacing \(\theta_i\) with \(\theta_i^\diamond\). Next, we want to compare \(\theta_i^\diamond\) with \(\theta_i^L\).
\[m(\theta^\diamond) = m'(0)\theta^\diamond + R, \qquad R = O((\theta^\diamond)^2)\]
Let \(m'(0) = 1\) WLOG. We have
\[\theta_i^\diamond = \theta_{i-1}^\diamond - \gamma(\theta_{i-1}^\diamond + R_{i-1}) + \gamma \eta_i\]

\[\theta_i^L = \theta_{i-1}^L - \gamma \theta_{i-1}^L + \gamma \eta_i\]

Let \(M_i = \theta_i^\diamond - \theta_i^L\). Then
\[M_i = (1 - \gamma)M_{i-1} - \gamma R_{i-1}\]
We can view this as an AR(1) process with innovations \(R_i\) that are dependent. Unraveling the recursion, we get

\textbf{TODO}

It suffices to show that \(\|R_i\|_2 = o(\sqrt \gamma)\).

We have
\[|R| \leq C |\theta^\diamond|^2\]

\[\E[R^2] \leq C^2 \E[|\theta^\diamond|^4] = \gamma^2\]

From the fact that \(\|\theta^\diamond\|_4^2 = O(\gamma)\). An interesting trick we can do to see this is
\[\min(|\theta|, |\theta|^2) \leq |\theta|^{p/2}\]
when \(2 < p < 4\). We can see this by breaking up into casework depending on whether \(|\theta| > 1\).

\begin{center}\rule{0.5\linewidth}{0.5pt}\end{center}

We can also prove a tail probability bound to control
\[P(|\theta_i| \geq y)\]
We can use the \textbf{Nagaev inequality}. Suppose we have \(\xi_1, \xi_2, \dots\) independent and mean \(0\), and that \(p \geq 2\). Set
\[\mu_{i,p} = \E[|\xi_i|^p]\]
Then
\[P(|\xi_1 + \cdots + \xi_n| \geq y) \leq C_1 \frac{\sum E[\xi_i|^p}{y^p} + C_2 \exp(-C_3 y^2 / \sum \E|\xi_i|^2)\]

\begin{center}\rule{0.5\linewidth}{0.5pt}\end{center}

SGD with decreasing \(\gamma_i\)
\[\theta_i = \theta_{i-1} - \gamma_i g(\theta_{i-1}, X_i)\]
Choosing \(\gamma\) is still a fundamental issue in this theory. Some example choices are \(\gamma_i = C i^{-\alpha}\) for \(\alpha \geq 0\). Other choices are cosine decay and annealed schedules.

In general, we want to choose the learning rate such that
\[\sum_i \gamma_i = \infty, \quad \sum_i \gamma_i^2 < \infty\]
For the constant learning rate case, we proved that
\[\|\theta_n\|_p^2 = O(\gamma)\]
In general, we have
\[\|\theta_n\|_p^2 = O(\gamma_n)\]
For small \(\gamma\),
\[(1 - \gamma_1) \cdots (1 - \gamma_k) \approx e^{-\gamma_1 - \cdots - \gamma_k} \approx e^{-c k^{1 - \alpha}}\]
This shows convergence to \(0\) at a sub-exponential rate. Indeed, we can use this to prove
\[\frac{\theta_n}{\sqrt \gamma_n} \Rightarrow \N(0, \Sigma)\]
Thus, we call this a last-iterate CLT. In 1992, Polyak and Juditsky proved that averaging can improve that rate of performance. If
\[\overline \theta_n = \frac{\theta_1 + \cdots + \theta_n}{n}, \qquad \|\overline \theta_n\| = O(1/\sqrt n)\]
Then
\[\sqrt n \hat \theta_n \Rightarrow \N(0, \Sigma)\]
We call this ASGD (averagedd SGD).

\textbf{Homework}: Let \(\theta_i = \theta_{i-1} - \gamma (\theta_{i-1} + X_i)\) where \(\gamma_i = C i^{-\alpha}\) where \(0 < \alpha < 1\). Then
\[\overline \theta_n = \sum_{i=1}^n \omega_{n,i} X_i\]
For some weights \(\omega_{n,i}\). Prove that \(\omega_{n,i} \approx 1/n\). Thus, we can approximate \(\overline \theta_n\) with \(\sum_{i=1}^n X_i\). In fact, show that
\[\sum_{i=1}^n (\omega_{n,i} - 1/n)^2 \rightarrow 0\]
as \(n \rightarrow \infty\). Furthermore, Ruppert in 1989 also introduced the technique of averaging in a technical report.

Following the introduction of ASGD, people proposed that more weight should be given to the later parameter values in the moving average definition of \(\overline \theta_n\). For example, an exponentially weighted moving average. In general, let us consider
\[\sum_{i=1}^n \theta_i \omega_{n,i}, \qquad \sum_{i=1}^n \omega_{n,i}\]
How can we optimally choose the weigts? There is a nice picture which shows this. Check the paper by Ziyang, Zhu, Wu (AISTATS 2026).
